\section{Classes, currents, and exceptional divisors}\label{sec:analytic}

The scaling argument transports one K\"ahler form from the original
space through several birational models.  Its trace need not be smooth,
or even have local potentials without further argument.  We establish
the precise trace relation, its positivity, and the comparison rules
used below.  Equalities between Bott--Chern classes will be distinguished
from equalities between divisors or currents.

\subsection{Pullback of classes}

We use the following cohomological descent theorem.  A space is in
\emph{Fujiki's class $\mathcal C$} if it is bimeromorphic to a compact
K\"ahler manifold.

\begin{lemma}[Pullback and descent]\label{lem:pullback-image}
Let $f:T\to Z$ be a proper surjective morphism with connected fibres between normal
compact analytic spaces with rational singularities.  Suppose that $T$
is in class $\mathcal C$ and that either the general fibre is rationally
connected or $R^jf_*\cO_T=0$ for every $j>0$.  Then
\[
 f^*: \BC(Z)\longrightarrow\BC(T)
\]
is injective, and its image consists exactly of the classes having
degree zero on every $f$-vertical curve.  If $T$ and $Z$ are K\"ahler,
a class on $Z$ is nef if and only if its pullback is nef.
\end{lemma}

\begin{proof}
The pullback assertion is \cite[Lemma~2.39]{HX2026}; nef descent is
\cite[Lemma~2.38]{DHP2024}.  In particular, the first assertion applies
to any proper bimeromorphic morphism between compact spaces with
rational singularities whose source is in class $\mathcal C$: its
general fibre is a point.  For a projective klt log Fano contraction,
relative vanishing supplies the higher-direct-image hypothesis once
rationality of the base has been established.
\end{proof}

\subsection{Positive currents on a normal model}

An $f$-exceptional real divisor is a finite real linear combination of
prime divisors whose images have codimension at least two.
We next give the descent statement needed when an exceptional
correction has either sign.  Pseudo-effectivity requires a positive
current with local plurisubharmonic potentials; we do not assume
that an arbitrary positive current on a singular space has them.

\begin{lemma}[Exceptional descent]\label{lem:positive-descent}
Let $p:U\to T$ be a resolution of a normal compact K\"ahler space,
with $U$ compact K\"ahler.  Let $\alpha\in\BC(T)$ and let $E$ be a
$p$-exceptional real divisor.  If $p^*\alpha+[E]$ is pseudo-effective,
then $\alpha$ is pseudo-effective.  No sign condition on $E$ is needed.
Pseudo-effective classes pull back to resolutions, and nef classes
on $T$ are pseudo-effective.
\end{lemma}

\begin{proof}
Choose a smooth representative $\theta$ of $\alpha$ with local smooth
potentials.  On the smooth compact K\"ahler space $U$, a positive
representative of the given class has the form
\begin{equation}\label{eq:positive-upstairs}
 S=p^*\theta+[E]+\ddc v
\end{equation}
for a global distribution $v$.  Remove from $T$ its singular locus
and the image of the exceptional locus, obtaining a smooth open set
$T^\circ$ whose complement has codimension at least two.  The map
$p$ is an isomorphism there.  Equation~\eqref{eq:positive-upstairs}
gives a global $\theta$-plurisubharmonic function $v^\circ$ on
$T^\circ$.  We do not assert that $v$ is quasi-plurisubharmonic along
$E$.

On a coordinate neighbourhood where $\theta=\ddc\rho$, extend
$\rho+v^\circ$ first across the removed codimension-two subset of
the regular locus, and then across the singular locus.  The first
extension is the plurisubharmonic Hartogs theorem; the second is its
normal-space version of Grauert--Remmert, recalled in
\cite[Section~2.1, pp.~927--928]{CMM2017}.  These extensions are unique.
On overlaps their differences are the original smooth
pluriharmonic differences of the functions $\rho$.  Subtracting
$\rho$ therefore gives a global potential $v_T$, and
\[
 \theta+\ddc v_T\geq0
\]
is a current with local plurisubharmonic potentials representing
exactly $\alpha$.  This proves descent, including for signed $E$.

For pullback, compose local plurisubharmonic potentials with the
resolution.  They cannot become identically $-\infty$ on a nonempty
open subset, since a resolution is an isomorphism on a dense open
set.  Thus they define the pulled-back positive current; see
\cite[Section~2.1, p.~928]{CMM2017}.

Finally, if $\alpha$ is nef, then $p^*\alpha$ is nef on $U$.
For a K\"ahler form $\omega_U$, choose positive representatives of
$p^*\alpha+\varepsilon[\omega_U]$.  Their masses are bounded as
$\varepsilon\downarrow0$, so a weak limit is a positive current in
$p^*\alpha$.  On the smooth K\"ahler manifold $U$ it has local
plurisubharmonic potentials.  The first part, with $E=0$, descends
this current to the required class on $T$.
\end{proof}

\subsection{Generalized pairs and negativity}

We recall the generalized-pair language in the form needed here
\cite[Section~2.1 and Definition~2.7]{DHY2023}.  A closed \emph{b-$(1,1)$
current} is a compatible collection of closed $(1,1)$-currents on
proper birational models: pushforward along a morphism between models
recovers the current on the lower model.  It \emph{descends} to a
model $U$ if its trace there and all higher traces have local
potentials, and the classes of the higher traces are pullbacks of
its class on $U$.
It is \emph{b-nef} if that class is nef on some such model.
We also use the relative version for real combinations of line-bundle
data on a carrier projective over a specified base.  Here nefness
means nonnegative degree on curves contracted to that base.  The
structure equation and discrepancy definition below are unchanged.

A generalized pair on $T$ consists of a boundary $B\geq0$, a b-nef
datum $\mathbf M$, and a log resolution $\nu:U\to T$ on which the
datum descends, together with an actual real divisor $B_U$ such that
$\nu_*B_U=B$ and
\begin{equation}\label{eq:generalized-structure}
 [K_U+B_U]+[\mathbf M_U]=\nu^*L,
 \qquad L\in\BC(T).
\end{equation}
The support of $B_U$ may be taken to have simple normal crossings.
Negativity makes this structure boundary unique; passing to higher
resolutions gives the other structure boundaries.  The generalized
log discrepancy of a prime $P$ on such a resolution is
\[
 a(P,T,B+\mathbf M)=1-\operatorname{coeff}_P B_U.
\]
The pair is generalized klt, or \emph{gklt}, if all these discrepancies
are positive, and generalized log canonical, or \emph{glc}, if they
are all nonnegative.  The \emph{gklt locus} is the open subset where
all discrepancies of primes centred there are positive; its
complement is the generalized non-klt locus.  We use the word
\emph{terminal} for an ordinary or
generalized klt pair whose discrepancies at all exceptional primes
are strictly greater than one.  Negative coefficients are permitted
in $B_U$, although the boundary $B$ on the model is effective.
Generalized klt spaces have rational singularities
\cite[Theorem~2.19]{DHY2023}.

The negativity lemma concerns actual divisors, not arbitrary
Bott--Chern classes.  We record also the strict form needed for
fourfold termination.

\begin{lemma}[Negativity and strict comparison]\label{lem:negativity}
\leavevmode
\begin{enumerate}[label=\textup{(\roman*)}]
\item Let $f:T\to Z$ be a proper bimeromorphic morphism of normal
complex spaces and let $E$ be a real Cartier divisor such that $-E$
is $f$-nef.  Then $E\geq0$ if and only if $f_*E\geq0$.
Consequently an $f$-exceptional, $f$-numerically trivial real Cartier
divisor is zero.
\item If, in addition, $f$ has connected projective fibres and
$E\geq0$ is $f$-anti-nef, then every fibre is either contained in
$\Supp E$ or disjoint from $\Supp E$.
\item Consider a nontrivial small birational diagram
\[
 T\xrightarrow{\ f\ }Z\xleftarrow{\ f^+\ }T^+
\]
of normal compact spaces, with projective contractions on both
sides.  Let $B\geq0$ be a real boundary, let $B^+$ be its strict
transform, and suppose the ordinary adjoints
$D=K_T+B$ and $D^+=K_{T^+}+B^+$ are real Cartier, with $-D$ relatively
ample on the left and $D^+$ relatively ample on the right.
The two exceptional loci have the same image $A\subset Z$.  On a
common resolution $p:V\to T$, $q:V\to T^+$,
\begin{equation}\label{eq:ordinary-flip-comparison}
 p^*D-q^*D^+=F,\qquad F\geq0,
\end{equation}
where $F$ is an actual divisor exceptional over both models.
For every prime divisor $P$ over these spaces,
\[
 a(P,T^+,B^+)\geq a(P,T,B),
\]
and the inequality is strict if the centre of $P$ on either side
is contained in that side's exceptional locus.
\item For the diagram in \textup{(iii)}, if $\dim T=4$, then $\dim A\leq1$ and
\[
 \dim\Exc(f)+\dim\Exc(f^+)\geq3.
\]
\item The discrepancy conclusions in \textup{(iii)} also hold for
generalized pairs with the same fixed b-nef datum and transformed
boundaries in a projective small diagram, provided the source log
class has negative degree on every curve contracted on that side,
and the target log class has positive degree on every curve
contracted on the other side.  They hold likewise for a projective divisorial
contraction $f:T\to T'$ when the target boundary is the pushforward,
the b-datum is unchanged, and the source log class has negative
degree on every $f$-vertical curve.  In this divisorial case strictness holds at every prime
whose source centre lies in $\Exc(f)$.
This part also applies when the fixed b-datum is nef only over a
specified base.
\end{enumerate}
\end{lemma}

\begin{proof}
Part~(i) is the analytic negativity lemma
\cite[Lemma~2.5]{HX2026}; apply it to both signs for the last assertion.
For (ii), points of a connected projective fibre can be joined by a
chain of irreducible curves.  If the fibre met both $\Supp E$ and
its complement, some curve in such a chain would meet the support
without being contained in it.  Its intersection with the effective
real Cartier divisor would be strictly positive, contrary to
$f$-anti-nefness.  The assertion for a real Cartier divisor follows
locally from positive combinations of effective rational Cartier
divisors, as explained below.

For (iii), let $G$ be the normalization of the graph in
$T\times_ZT^+$, with projections $p_G,q_G$ and map $h:G\to Z$.
Compatible canonical data give the actual real Cartier divisor
\[
 F_G=p_G^*D-q_G^*D^+.
\]
This difference is defined locally by compatible meromorphic canonical
generators and therefore glues, even when the canonical sheaves were
specified as $\Q$-line bundles.  Smallness makes it exceptional over
both sides.  It is effective by (i), applied over $T$: on a
$p_G$-vertical curve the degree of $-F_G$ is the degree of
$q_G^*D^+$, which is nonnegative.

For every $h$-vertical curve $C$ one has
\begin{equation}\label{eq:graph-strict-degree}
 F_G\cdot C<0.
\end{equation}
Indeed, normalization is finite onto the graph, so the two projections
cannot both contract $C$.  The nonzero degrees contributed by $p_G^*D$
and $-q_G^*D^+$ are both negative.  An effective real Cartier divisor
has nonnegative degree on a curve outside its support.  For completeness,
this fact and its strict pullback analogue hold even if its individual
prime components are not $\Q$-Cartier: locally express it in the real
span of finitely many Cartier divisors.  Vanishing of coefficients
outside its support and nonnegativity of the remaining coefficients
are rational linear conditions.  The resulting rational polyhedral
cone expresses it as a positive real combination of effective
$\Q$-Cartier divisors.  The usual effective Cartier assertions apply
to each term; summing the nonnegative local intersection numbers
gives the asserted global degree inequality.  In particular, the pullback has positive coefficient
at every prime whose centre is contained in the support.

The morphism $p_G$ has connected fibres because $T$ is normal;
hence $h=f p_G$ has connected fibres.  Its fibres are projective
because $G$ is finite over the graph in the fibre product.  Each
point of a positive-dimensional fibre lies on a curve in that fibre.  Equation
\eqref{eq:graph-strict-degree} therefore places the entire fibre in
$\Supp F_G$.  To identify the relevant fibres, suppose $f$ is an
isomorphism over an open subset of $Z$.  There the log divisor on
$T^+$ is pulled back from $Z$, by smallness and agreement in codimension
one.  Relative ampleness excludes an $f^+$-vertical curve there;
projectivity, connectedness and normality then make $f^+$ an
isomorphism there as well.  Interchanging the sides proves that their
exceptional images coincide.  Thus $h^{-1}(A)\subset\Supp F_G$.

Pull $F_G$ back to a resolution on which $P$ appears.  Its coefficient
at $P$ is exactly
$a(P,T^+,B^+)-a(P,T,B)$, with our log-discrepancy convention.
Effectivity gives monotonicity, and the support assertion gives the
claimed strictness.  This proves (iii).

For (iv), a small exceptional locus in a fourfold has dimension at
most two and has positive-dimensional fibres, so its image has
dimension at most one.  Moreover $F_G\ne0$ by
\eqref{eq:graph-strict-degree}.  A component of its support has
dimension three and maps finitely into
$\Exc(f)\times_Z\Exc(f^+)$.  The dimension of this product is at most
the sum of the dimensions of its two factors.  This proves (iv).

For (v), choose a common smooth resolution $p:V\to T$,
$q:V\to T'$ carrying the fixed b-datum and projective over the
common contraction base.  Such a resolution is obtained by taking
projective modifications of the graph that resolve the maps to a
carrier of the datum.  Subtract the actual
structure boundaries to obtain a real Cartier divisor
$F=B_{T,V}-B_{T',V}$.  Its coefficients are the differences of
generalized log discrepancies, and cancellation of the fixed
b-part gives
\[
 [F]=p^*L_T-q^*L_{T'}.
\]
Let $h:V\to Z$ be the morphism to the common contraction base;
in the divisorial case $Z=T'$ and $q=h$.  It is projective by this
choice and has connected fibres because it is bimeromorphic onto
the normal space $Z$.  The boundary-pushforward
condition makes $F$ exceptional over $Z$.  The relative signs of
the log classes give $F\cdot C\leq0$ on every $h$-vertical curve,
with strict inequality whenever $p(C)$ or, in the small case,
$q(C)$ is a curve.  Thus (i) gives $F\geq0$.

If $y$ lies in an exceptional image, a vertical curve downstairs
can be lifted to a curve on $V$ dominating it: use a component of
its projective inverse image and cut by relatively ample
hyperplanes.  Such a curve has negative $F$-degree and hence lies
in $\Supp F$.  Part~(ii) now puts the entire fibre $h^{-1}(y)$ in
the support.  Pullback to a higher resolution has positive
coefficient at every prime centred in this support.  This proves
all the strict discrepancy inequalities, including for a centre
on the flipped side.  No trace current on the normalized graph
is assumed to be Cartier; the actual divisor is constructed from
the structure boundaries on the smooth resolution.
\end{proof}

\subsection{A fixed nef b-part and its traces}

For the finite-program proof, fix a K\"ahler form $h$ on the original
space $X$, and set $H=[h]$.  It defines a positive b-current
$\mathbf H$: for a marked birational model $T$, choose a common
resolution $p:V\to T$, $q:V\to X$ and put
\begin{equation}\label{eq:actual-H-trace}
 \mathbf H_T=p_*q^*h.
\end{equation}
Compatibility on higher resolutions makes this independent of the
choice of resolution \cite[Claim~2.5]{DHY2023}.  Throughout the proof
this is one fixed current datum; we do not replace it by arbitrary
cohomologous currents.  The class of the trace, when it exists, will
be denoted by $H_T$.

\begin{lemma}[The exceptional trace relation]\label{lem:class-trace}
Let $T$ be a normal compact K\"ahler space with rational singularities,
marked birationally to $X$.  Suppose that on a common smooth K\"ahler
resolution $p:V\to T$, $q:V\to X$ there are a class $H_T\in\BC(T)$
and an actual $p$-exceptional real divisor $J_T$ satisfying
\begin{equation}\label{eq:class-trace}
 p^*H_T=q^*H+[J_T].
\end{equation}
Then $J_T\geq0$.  The class $H_T$ and the divisor $J_T$ are unique,
and the correction on a higher resolution is the pullback of $J_T$.

For an effective boundary $B_T$ with real Cartier ordinary adjoint,
the generalized structure boundary for the nef datum $t\mathbf H$ is
\begin{equation}\label{eq:trace-boundary}
 B_{T,V}(t)=p^*(K_T+B_T)-K_V+tJ_T,\qquad t\geq0.
\end{equation}
In particular, decreasing $t$ increases every generalized log
discrepancy.  If the generalized pair is gklt, its ordinary pair is
klt; if all its exceptional log discrepancies exceed one, the
ordinary pair is terminal.

Conversely, suppose $t>0$ and a generalized pair with boundary $B_T$
and this fixed datum $t\mathbf H$ is given by
\eqref{eq:generalized-structure}, with log class $L$ and real Cartier
ordinary adjoint.  Then \eqref{eq:class-trace} holds with
\[
 H_T=\frac{L-[K_T+B_T]}{t}.
\]
\end{lemma}

\begin{proof}
On every $p$-vertical curve, $-J_T$ has the same degree as the nef
class $q^*H$.  Lemma~\ref{lem:negativity}(i) gives $J_T\geq0$.
Subtract two possible relations.  The difference of the correction
divisors is $p$-exceptional and numerically trivial over $T$, hence
zero.  Pullback injectivity from Lemma~\ref{lem:pullback-image} then
gives uniqueness of the class as well.

Pulling back \eqref{eq:class-trace} to a higher resolution gives the
claimed correction there by uniqueness.  Conversely, if the relation
is initially known on a higher resolution $r:V'\to V$, push its
correction down to $V$.  The difference between the upstairs
correction and the pullback of this pushforward is $r$-exceptional
and numerically trivial over $V$, so is zero.  Injectivity of $r^*$
gives the relation on $V$.

Equation~\eqref{eq:trace-boundary} follows by adding $K_V$ and the
fixed nef class $tq^*H$ to its two sides.  Its coefficients give
\[
 a(P,T,B_T+t\mathbf H)
 =a(P,T,B_T)-t\operatorname{coeff}_P J_T
\]
on every sufficiently high resolution.  This proves the discrepancy
assertions.

For the converse, let $B_V$ be the given actual generalized structure
boundary and put $B_V^{\rm ord}=p^*(K_T+B_T)-K_V$.  Both push forward
to $B_T$, so
\[
 J_T=\frac{B_V-B_V^{\rm ord}}{t}
\]
is an actual $p$-exceptional divisor.  Subtract the ordinary log class
from \eqref{eq:generalized-structure} to obtain
\eqref{eq:class-trace}.  This constructs the divisor before applying
negativity; it does not infer a divisor from an arbitrary class
difference.
\end{proof}

The converse is particularly useful for an extraction carrying a
pulled-back generalized log class.  It shows that the divided
difference defining its trace is the same class whenever the marked
model is the same, even if its boundary was obtained at a different
scaling parameter.

\begin{corollary}[Positivity of the trace]\label{cor:big-trace}
Under the hypotheses of Lemma~\ref{lem:class-trace}, the actual trace
$\mathbf H_T$ has local plurisubharmonic potentials and represents
$H_T$.  The class $H_T$ is big.
\end{corollary}

\begin{proof}
The positive current $q^*h+[J_T]$ represents $p^*H_T$.
The proof of Lemma~\ref{lem:positive-descent} constructs a positive
current $S_T$ in $H_T$ with local potentials.  On the big open where
$p$ is an isomorphism and the correction is absent, it agrees with
$p_*q^*h$.  Their difference is a closed current of order zero
supported in codimension at least two.  Such a current vanishes:
after a local embedding it has bidimension
$(\dim T-1,\dim T-1)$, whereas its support has smaller dimension
\cite[Chapter~III, Corollary~2.11]{DemaillyCADG2012}.
Thus $S_T$ is precisely the actual trace in
\eqref{eq:actual-H-trace}.  This step is needed because a trace of
a positive b-current need not have local potentials a priori
\cite[Remark~2.6(ii)]{DHY2023}.

The class $q^*H$ is nef and has positive top self-intersection:
$q^*h$ is semipositive and is positive on the dense open where $q$
is an isomorphism.  By \cite[Theorem~0.5]{DemaillyPaun2004}, it contains
a K\"ahler current $R\geq\delta\omega_V$ for some $\delta>0$ and
some K\"ahler form $\omega_V$.  Fix a K\"ahler form $\omega_T$ and
choose $C>0$ with $p^*\omega_T\leq C\omega_V$.  For
$0<\varepsilon<\delta/C$,
\[
 R+[J_T]-\varepsilon p^*\omega_T\geq0.
\]
Its class is $p^*(H_T-\varepsilon[\omega_T])$.
Lemma~\ref{lem:positive-descent} shows that
$H_T-\varepsilon[\omega_T]$ is pseudo-effective.  Hence $H_T$
contains a K\"ahler current, as required.
\end{proof}

\subsection{Comparison with the same nef datum}

The preceding results supply the positivity needed for scaling.  We
finish with the comparison that later excludes divisorial steps in
the limiting program and identifies the two nef pullbacks in the
final contradiction.

\begin{lemma}[Small-model comparison]\label{lem:small-comparison}
Let $(T,B_T+\mathbf M)$ and $(T',B_{T'}+\mathbf M)$ be generalized
pairs with the same b-nef datum, whose log classes are $L_T$ and
$L_{T'}$.  Suppose the marked birational map $T\dashrightarrow T'$
extracts no divisors and $B_{T'}$ is the pushforward of $B_T$.
On a common sufficiently high resolution $p:V\to T$, $q:V\to T'$, put
\[
 F=B_{T,V}-B_{T',V}.
\]
Then $F$ is an actual divisor exceptional over $T'$, and
\begin{equation}\label{eq:small-comparison}
 p^*L_T-q^*L_{T'}=[F],\qquad
 \operatorname{coeff}_P F
 =a(P,T',B_{T'}+\mathbf M)-a(P,T,B_T+\mathbf M).
\end{equation}
If the map is small, $F$ is exceptional over both models.  In this
case nefness of $L_{T'}$ implies $F\geq0$, and nefness of both log
classes implies $F=0$.  These conclusions also hold relatively over
a common base.  In particular equal pullback classes imply equality
of the structure boundaries.
\end{lemma}

\begin{proof}
Subtract the two structure equations on $V$.  The trace of the fixed
b-datum cancels, giving \eqref{eq:small-comparison}.  At a prime
dominating a divisor on $T'$, the corresponding boundary coefficient
on $T$ is unchanged, so its coefficient in $F$ is zero.  This proves
exceptionality over $T'$.  Smallness gives the same assertion over
$T$.

If $L_{T'}$ is nef, then $-F$ is $p$-nef, because its degree on a
$p$-vertical curve is the degree of $q^*L_{T'}$.  Negativity gives
$F\geq0$.  If $L_T$ is also nef, then $F$ is $q$-nef, so applying
negativity with the opposite sign gives $F\leq0$.  The same proof
uses only relative nefness when the models lie over a common base.
For equal pullback classes, the already constructed exceptional
divisor is relatively numerically trivial, and is zero by
Lemma~\ref{lem:negativity}(i).
\end{proof}

All equalities of pulled-back log classes in this lemma are
cohomological.  Equality of actual currents was used only in the
separate trace-identification argument.  These distinctions allow
the ordinary steps and the generalized auxiliary constructions to
share one fixed nef datum.
